\documentclass[UTF8,a4paper,11pt]{ctexart}

\usepackage[left=2.15cm,right=2.15cm,top=1.85cm,bottom=1.9cm]{geometry}
\usepackage{booktabs,tabularx,array,xcolor,enumitem,hyperref}
\hypersetup{colorlinks=true,linkcolor=blue,urlcolor=blue}
\setlength{\parindent}{2em}
\setlength{\parskip}{0.32em}
\renewcommand{\arraystretch}{1.42}
\newcolumntype{Y}{>{\raggedright\arraybackslash}X}
\definecolor{accent}{HTML}{146C73}
\definecolor{soft}{HTML}{EAF6F5}
\definecolor{warn}{HTML}{FFF3D6}
\newcommand{\prompt}[1]{\textbf{“#1”}}

\title{化小北 AI 眼镜教师版操作指南}
\author{}
\date{}

\begin{document}

\maketitle
\vspace{-1.5em}
\begin{center}
  \small 本指南介绍眼镜的基础操作、教师版语音功能和教师网页工作台。
\end{center}

\section{基础操作}

眼镜主要通过右镜腿上的按钮和滑动手势操作：

\begin{description}[leftmargin=3.8em,style=nextline]
  \item[长按右镜腿按钮] 开机或关机。
  \item[单击右镜腿按钮] 进入当前选中的应用或内容，相当于“确认”键。
  \item[快速双击右镜腿按钮] 返回上一层，或退出当前应用。
  \item[向左或向右滑动] 移动选项、切换页面或控制当前内容。
\end{description}

开机后，在应用列表中左右滑动，找到\textbf{教师版 App}。选中后单击右镜腿按钮，即可进入教师版。

\section{进入软件后的操作}

\subsection{启动成功的表现}

进入教师版后，系统会自动连接校园网服务。连接成功时，眼镜会说：

\begin{center}
  \fcolorbox{accent}{soft}{\parbox{0.76\textwidth}{\centering\Large\color{accent}\prompt{助手就绪，请说指令。}}}
\end{center}

听到这句话，就表示教师版已经正常启动。提示音播放完毕后，系统会\textbf{自动开始持续录音并上传教师音频记录}，不需要再说“开始录音”。教师版以语音操作为主，画面保持简洁或没有长期显示内容属于正常现象。

\begin{center}
  \fcolorbox{orange!70!black}{warn}{\parbox{0.88\textwidth}{
  \textbf{请特别注意：}从听到“助手就绪，请说指令”开始，教师版会一直记录声音；拍照前、录像过程中以及停止录像后，音频记录都会继续。结束本次使用并退出教师版后，持续录音才会停止。
  }}
\end{center}

\subsection{怎样发出指令}

正常操作顺序如下：

\begin{enumerate}[leftmargin=2.4em]
  \item 先说唤醒词 \prompt{小北助手}。
  \item 稍等几秒，听到系统回答 \prompt{我在}。
  \item 紧接着清楚地说出“拍照”或“开始录像”等指令。
\end{enumerate}

听到“我在”表示唤醒成功。系统会留出约三秒钟接收指令，建议一次只说一件事。如果听到 \prompt{我没有听清，请再说一遍}，说明本次没有识别成功；等系统重新进入监听状态后，再从“小北助手”开始说一遍。

\subsection{拍照}

推荐说法：

\begin{itemize}[leftmargin=2.4em]
  \item “小北助手”——听到“我在”后说“拍张照片”；
  \item 或说“把现在这个画面拍下来”“帮我拍一下”。
\end{itemize}

识别成功后，眼镜会说 \prompt{开始拍照，请保持画面稳定}。听到这句话后，请短暂保持头部和视野稳定；系统会完成拍照并自动上传，然后恢复监听。

\subsection{开始录像}

需要录像时，请明确说出“录像”或“视频”，例如：

\begin{itemize}[leftmargin=2.4em]
  \item “小北助手”——听到“我在”后说“开始录像”；
  \item “把这一段实验过程录成视频”；
  \item “开始录制视频”。
\end{itemize}

识别成功后，眼镜会说 \prompt{开始记录}。\textbf{在教师版中，听到这句提示就表示相机录像已经开始。}与此同时，原有的持续音频记录仍会进行。

为了避免混淆，教师需要录像时建议直接说“开始录像”；教师版的音频记录本来就会自动持续进行，不需要另发“开始录音”指令。

\subsection{停止录像}

录像过程中，\textbf{不需要再次说“小北助手”}，直接说下面任意一句即可：

\begin{itemize}[leftmargin=2.4em]
  \item “停止录像”或“结束录像”；
  \item “保存录像”；
  \item “录好了”“可以了”“先停一下”。
\end{itemize}

系统识别后会结束并上传这段录像，然后说 \prompt{记录已停止}。听到这句话，表示本段视频已经结束，系统已回到普通监听和持续录音状态。

\subsection{提示音速查}

{\small
\begin{tabularx}{\textwidth}{p{0.28\textwidth}p{0.31\textwidth}Y}
\toprule
当前情况 & 会听到 & 表示什么 \\
\midrule
进入教师版 & “助手就绪，请说指令” & 连接成功，持续录音即将开始 \\
说出“小北助手” & “我在” & 唤醒成功，请马上说具体指令 \\
拍照指令识别成功 & “开始拍照，请保持画面稳定” & 正在拍照，请保持视野稳定 \\
录像指令识别成功 & “开始记录” & 相机录像已经开始，音频也在继续记录 \\
停止录像识别成功 & “记录已停止” & 本段录像结束，持续录音仍在进行 \\
指令没有识别成功 & “我没有听清，请再说一遍” & 本次未执行动作，请重新唤醒并说明 \\
\bottomrule
\end{tabularx}
}

\subsection{结束本次使用}

使用结束时，按返回键或双击镜腿按键退出教师版。退出后持续录音停止；如果退出时仍在录像，系统会先尽量保存当前视频。

\section{教师网页工作台}

\subsection{登录与进入工作台}

在校园网环境下，用浏览器打开：

\begin{center}
  \url{http://162.105.209.134/login}
\end{center}

点击“进入北大门户登录”，完成北大账号认证后会直接进入教师工作台。系统依据门户返回的教师或职工身份自动显示教师功能，不需要手动选择身份。若页面误显示为学生工作台，请联系管理员调整。

教师的认领、资料查看、课程整理、学生数据、演示视频和意见反馈都在同一个工作台中，通过页面顶部标签切换，不需要打开其他后台地址。

\subsection{下课后认领本次数据}

眼镜不会永久绑定给某位教师。每次使用结束后，请根据眼镜外壳上的 01--09 编号和实际上课时间认领：

\begin{enumerate}[leftmargin=2.4em]
  \item 打开“数据认领”，先查看认领日历或“等待认领数据”；
  \item 选择眼镜编号，填写开始日期、时间和结束日期、时间；
  \item 点击“预览匹配数据”，检查数量和实验信息；
  \item 确认无误后点击“确认认领”。
\end{enumerate}

日历中会直观区分待认领和自己已经认领的数据。点击某一天的条目可展开具体时段；“最近一小时”和“结束 = 开始 + 1 小时”可用于快速填写。若眼镜编号显示未校准，请联系管理员，不要选择其他编号代替。

\subsection{查看自己的音频、录像和照片}

“我的数据”汇总本账号已经认领的资料；“教师资料”进一步按音频、实验录像和照片分类。刚上传的内容可能需要稍等片刻，再点击“刷新”。

教师只能查看、下载或删除自己认领的教师资料。还没有认领的数据不会出现在个人资料中；其他教师认领的资料也不会显示。

\subsection{生成时间段课程报告}

打开“课程整理”，按下面顺序操作：

\begin{enumerate}[leftmargin=2.4em]
  \item 选择“课程内容整理”或“共性问题与指导建议”；
  \item 选择开始日期和结束日期，可覆盖当天，最长 31 天；
  \item 根据需要选择实验和眼镜编号；
  \item 点击“查找可用记录”，检查并勾选要使用的具体记录；
  \item 点击“生成报告”，等待页面显示完成后查看或下载 Markdown。
\end{enumerate}

报告只会使用\textbf{当前教师已经认领且完成文字转写的教师音频}。照片、纯视频、待认领资料、其他教师资料以及尚未完成转写的录音不会进入报告。课程内容整理至少需要一条有效转写；共性问题报告必须选择一个实验，并至少选择两条不同记录。

报告生成过程中会显示排队或处理中。生成失败或资料不足时，页面会直接说明原因；调整日期、实验、眼镜或勾选记录后可以重新生成。旧的单日课程报告仍可在历史列表中查看。

\subsection{学生数据、演示视频和意见反馈}

“学生数据”用于查看学生眼镜上传的资料。“演示与反馈”中可以为指定实验上传 MP4 演示视频，也可以填写教学问题和改进建议。上传演示视频前先确认实验类型正确。

\subsection{常见情况}

{\small
\begin{tabularx}{\textwidth}{p{0.34\textwidth}Y}
\toprule
现象 & 处理办法 \\
\midrule
进入教师版后没有听到就绪提示 & 确认眼镜已连接校园网，退出软件后重新进入；仍无提示时联系管理员。 \\
说“小北助手”后没有回答 & 稍等几秒；面向前方、适当提高音量后再说一次。 \\
反复听到“我没有听清” & 听到“我在”后再说完整短句，一次只说拍照或录像中的一项。 \\
网页显示为学生工作台 & 联系管理员把门户身份调整为教师；以后再次登录会保留管理员设置。 \\
找不到刚上传的资料 & 先等待上传完成，再按实际眼镜编号和时间认领，然后刷新“教师资料”。 \\
记录显示“转写中” & 暂时不能用于报告，稍后刷新并等待状态变为可用。 \\
共性问题无法生成 & 选择一个实验，并确认至少勾选两条已完成转写的教师音频。 \\
眼镜编号显示不正确 & 记下眼镜外壳编号并联系管理员校准，不要自行改用其他编号。 \\
\bottomrule
\end{tabularx}
}

\end{document}
